\section{Evaluasi}
\label{sec:evaluasi-bab14}

\subsection{Soal Pilihan Ganda}

\begin{enumerate}
    \item Setelah \asm{PUSH AX} dijalankan, nilai \reg{SP} \dots
          \begin{enumerate}
              \item bertambah 2
              \item berkurang 2
              \item tidak berubah
              \item bertambah 4
          \end{enumerate}

    \item Alamat kembali pada pemanggilan prosedur disimpan oleh
          instruksi \dots
          \begin{enumerate}
              \item \asm{JMP} \item \asm{CALL}
              \item \asm{RET} \item \asm{PUSH}
          \end{enumerate}

    \item Setelah \asm{PUSH BP} dan \asm{MOV BP, SP}, parameter
          pertama pada stack frame diakses melalui \dots
          \begin{enumerate}
              \item \code{[BP+0]} \item \code{[BP+2]}
              \item \code{[BP+4]} \item \code{[BP+6]}
          \end{enumerate}

    \item Macro diakhiri dengan direktif \dots
          \begin{enumerate}
              \item \asm{ENDP} \item \asm{ENDM}
              \item \asm{ENDS} \item \asm{END}
          \end{enumerate}

    \item Perbedaan utama prosedur dan macro adalah \dots
          \begin{enumerate}
              \item prosedur memakai \asm{CALL}, macro disalin saat
                    perakitan
              \item prosedur tidak dapat memakai register
              \item macro tidak dapat memakai parameter
              \item macro selalu lebih kecil
          \end{enumerate}
\end{enumerate}

\subsection{Soal Uraian}

\begin{enumerate}
    \item Jelaskan pengertian stack dan sifat LIFO. Telusuri nilai
          \reg{SP} serta isi stack untuk \lstasm|PUSH AX| dan
          \lstasm|POP BX| bila \reg{SP} mula-mula \hex{FFFE}.

    \item Jelaskan fungsi \asm{PROC}, \asm{ENDP}, \asm{CALL}, dan
          \asm{RET}. Mengapa \asm{CALL} berbeda dari \asm{JMP}?

    \item Jelaskan cara mengakses parameter melalui \eng{stack frame}
          dengan \reg{BP}. Sebutkan isi \code{[BP+0]}, \code{[BP+2]},
          \code{[BP+4]}, dan \code{[BP+6]}, lalu jelaskan fungsi
          \asm{RET} \code{4}.

    \item Jelaskan cara mendefinisikan macro berparameter dan apa
          yang terjadi saat macro dipakai. Sebutkan empat perbedaan
          prosedur dan macro.
\end{enumerate}

\subsection{Kisi-kisi Penilaian}

\begin{table}[H]
    \centering
    \caption{Kisi-kisi evaluasi Bab~\ref{chap:prosedur}.}
    \label{tab:kisi-bab14}
    \small
    \begin{tabular}{@{}p{1.5cm}p{5.0cm}p{3.0cm}c@{}}
        \toprule
        \textbf{Sub-CPMK} & \textbf{Indikator yang diukur} &
        \textbf{Bentuk soal} & \textbf{Bobot} \\
        \midrule
        3.5 & Menjelaskan stack dan PUSH/POP &
        Pilihan ganda 1; Uraian 1 & 25\% \\
        \addlinespace
        3.5 & Menjelaskan prosedur, CALL, dan RET &
        Pilihan ganda 2; Uraian 2 & 25\% \\
        \addlinespace
        3.5 & Mengakses parameter melalui stack frame &
        Pilihan ganda 3; Uraian 3 & 25\% \\
        \addlinespace
        3.5 & Mendefinisikan macro dan membedakannya &
        Pilihan ganda 4, 5; Uraian 4 & 25\% \\
        \bottomrule
    \end{tabular}
\end{table}

\begin{checklistbox}
\vspace{1.1em}
\textbf{Checklist Pencapaian Kompetensi
Bab~\ref{chap:prosedur}}

Seorang mahasiswa dinyatakan kompeten pada Sub-CPMK 3.5 apabila:
\begin{itemize}
    \item[\checked] dapat menjelaskan sifat LIFO pada stack;
    \item[\checked] dapat menggunakan PUSH dan POP;
    \item[\checked] dapat menjelaskan perubahan SP;
    \item[\checked] dapat menulis prosedur dengan PROC/ENDP;
    \item[\checked] dapat menggunakan CALL dan RET;
    \item[\checked] dapat mengakses parameter melalui BP;
    \item[\checked] dapat membuat macro berparameter;
    \item[\checked] dapat membedakan prosedur dan macro.
\end{itemize}
\end{checklistbox}
